\section{Updated Problem and Objectives}

\subsection*{Updated problem statement}
This project investigates unsupervised industrial surface anomaly detection on the MVTec 3D-AD dataset \cite{bergmann2022mvtec3d}. The task is not formulated as standard supervised defect classification. Instead, models are trained only on defect-free samples and must assign high anomaly scores to abnormal test samples and, when possible, localize the affected regions. This setting follows the evaluation logic of MVTec 3D-AD and related industrial anomaly detection benchmarks \cite{bergmann2019mvtec,bergmann2022mvtec3d}.

The initial project plan compared explicit feature engineering against learned representations. During implementation, this formulation was refined. The first handcrafted baseline compressed each local patch into nine statistics, including mean height, standard deviation, gradient magnitude and roughness proxies. Qualitative inspection showed that this representation discarded too much spatial structure: small defects could become indistinguishable from sensor noise, and larger defects often required information beyond a single local descriptor. Therefore, the current classical baseline uses the raw normalized depth values of each patch, followed by PCA and a category-specific One-Class SVM \cite{scholkopf2001estimating}. This still remains a classical one-class model, but it is more directly comparable to the planned autoencoder because both methods operate on the same local depth-patch representation.

The current problem statement is thus more precise: we study how far local depth-only patch normality modeling can go on MVTec 3D-AD, and whether learned reconstruction models or additional RGB information are needed to capture defects that are not locally distinctive in geometry alone. This interim stage is especially useful because the first experiments already show that the choice of representation and context is as important as the choice of anomaly detector.
